\bookchapter{Capítulo 58: Guía de referencia rápida y recursos}

\chapterintro{Este capítulo es una guía de referencia rápida que puede usarse de forma independiente. Cuando no recuerdes alguna sintaxis o necesites decidir rápido, ven directo aquí. Trátalo como una “tarjeta de bolsillo”: código de consulta, preguntas frecuentes y recursos de aprendizaje, sin saltarte ni una página.}

\booksection{Tarjeta de referencia rápida de código}

\begin{center}
\begin{tabularx}{\linewidth}{|X|X|}
\hline
\textbf{Qué quieres hacer} & \textbf{Código clave} \tabularnewline
\hline
Guardar un estado mutable & var x by remember \{ mutableStateOf(valorInicial) \} \tabularnewline
Conservar el estado al girar la pantalla & rememberSaveable \{ mutableStateOf(valorInicial) \} \tabularnewline
Diseño en vertical u horizontal & \texttt{Column \{ \}}/\texttt{Row \{ \}} \tabularnewline
Lista con carga diferida & LazyColumn \{ items(lista) \{ ... \} \} \tabularnewline
Navegar a otra pantalla & navController.navigate("ruta") \tabularnewline
Tareas de espera dentro de una corrutina & \texttt{viewModelScope.\allowbreak launch \{ .\allowbreak .\allowbreak .\allowbreak  \}} \tabularnewline
Observar flujos de datos & \texttt{val list by flow.\allowbreak collectAsState()} \tabularnewline
Guardar ajustes pequeños & context.dataStore.edit \{ it[KEY] = valor \} \tabularnewline
Guardar datos estructurados & El trío de Room: @Entity / @Dao / @Database \tabularnewline
Descargar datos de la red & Interfaz Retrofit + funciones suspend \tabularnewline
Enviar una notificación & NotificationChannel + NotificationCompat.Builder \tabularnewline
Escribir logs & \texttt{Log.\allowbreak d("TAG", "message")} \tabularnewline
\hline
\end{tabularx}
\end{center}

\begin{center}\small\color{gray}Tarjeta de referencia rápida de desarrollo Android (edición de agosto de 2026)\end{center}

\booksection{Preguntas frecuentes (FAQ)}

P1: No sé nada de programación, ¿cuánto se tarda en aprender desarrollo Android?

Sigue el camino de este libro: en las semanas 1-2, instalar el entorno y aprender lo básico de Kotlin; en las 3-4, aprender Compose y construir interfaces; en las 5-6, terminar el proyecto de la app de gastos; en las 7-8, completar el proyecto del clima. Dos meses para pasar de cero a crear pequeñas Apps por tu cuenta. La condición: dedicar 1-2 horas al día y teclear código de verdad, no solo leer.

P2: ¿Todavía hace falta aprender Java?

No. Los proyectos nuevos, los tutoriales nuevos y los puestos nuevos van esencialmente con Kotlin. Solo al “mantener un proyecto de hace diez años” te toparás con mucho Java. Si aprendes Kotlin y luego el trabajo lo exige, ponerle Java será rápido: los dos viven en el mismo ecosistema de máquina virtual.

P3: ¿Emulador o dispositivo real?

Para desarrollar y depurar, el emulador es cómodo (instantáneas, varios modelos), pero la aceptación final debe hacerse en un dispositivo real: el tacto, el rendimiento, la red de verdad y los sensores no los reemplaza ningún emulador. Si puedes, comprar un móvil real de gama media como equipo de pruebas es la inversión más rentable.

P4: ¿Es fácil encontrar trabajo en desarrollo Android?

El mercado de 2026: los puestos junior que solo “saben hacer interfaces” tienen mucha competencia, pero los desarrolladores que “dominan Kotlin + conocen arquitectura + tienen apps publicadas” siguen siendo codiciados. Mientras aprendes, tu mejor currículum es esa App hecha por ti que está en la tienda.

P5: ¿Qué hacer cuando un error te bloquea?

Los tres pasos de oro: (1) busca en StackOverflow o en un buscador la primera línea del error tal cual; (2) consulta la página correspondiente de la documentación oficial; (3) si nada funciona, empaqueta “comportamiento esperado + lo que ocurre + el error completo + el código relacionado” y pregúntale a la IA o a la comunidad. La gran mayoría de los problemas se resuelven en el primer paso.

\booksection{Recursos de aprendizaje recomendados}

\begin{center}
\begin{tabularx}{\linewidth}{|X|X|X|}
\hline
\textbf{Tipo} & \textbf{Recurso} & \textbf{Por qué se recomienda} \tabularnewline
\hline
Oficial & developer.android.com & Documentación oficial de Android; incluye los tutoriales de Compose y Kotlin, la máxima autoridad \tabularnewline
Oficial & Sitio oficial de Kotlin (kotlinlang.org) & Documentación oficial del lenguaje + patio de pruebas en línea \tabularnewline
Oficial & Google Codelabs & Cursos prácticos paso a paso, amabilidad máxima con principiantes \tabularnewline
Comunidad & Stack Overflow (etiqueta android) & La primera parada para pedir ayuda con errores \tabularnewline
Comunidad & Sección Android de Juejin / CSDN & Muchos artículos prácticos en chino; ojo con distinguir nuevos de viejos (mira primero la fecha) \tabularnewline
Vídeo & Canal oficial de Google Developers & Las explicaciones de las novedades de cada conferencia anual de desarrolladores \tabularnewline
Complemento & Séptima parte, "Programación con IA en la práctica" (capítulos 37-47) & El método completo para desarrollar Android con ayuda de la IA \tabularnewline
\hline
\end{tabularx}
\end{center}

\begin{center}\small\color{gray}Lista de recursos para aprender desarrollo Android (agosto de 2026)\end{center}

\booksection{Comandos Gradle de referencia rápida}

En el día a día basta con pulsar botones, pero la línea de comandos a veces es más eficiente (sobre todo junto a herramientas de IA o pipelines de CI):

\begin{center}
\begin{tabularx}{\linewidth}{|X|X|}
\hline
\textbf{Comando} & \textbf{Función} \tabularnewline
\hline
\texttt{.\allowbreak /\allowbreak gradlew assembleDebug} & Generar un APK de depuración (en Windows, quita el ./) \tabularnewline
\texttt{.\allowbreak /\allowbreak gradlew assembleRelease} & Generar la versión de lanzamiento (requiere la firma configurada) \tabularnewline
\texttt{.\allowbreak /\allowbreak gradlew test} & Ejecutar todas las pruebas unitarias \tabularnewline
\texttt{.\allowbreak /\allowbreak gradlew lint} & Análisis estático: detecta bugs potenciales y malas prácticas \tabularnewline
\texttt{.\allowbreak /\allowbreak gradlew clean} & Limpiar los artefactos de compilación (la versión "reiniciar y arreglar" de la compilación) \tabularnewline
\texttt{.\allowbreak /\allowbreak gradlew :app: dependencies} & Ver el árbol de dependencias (el arma definitiva contra los conflictos de versiones) \tabularnewline
\hline
\end{tabularx}
\end{center}

\begin{center}\small\color{gray}Comandos Gradle de referencia rápida\end{center}

\begin{tcbinfo}{Consejo: si la compilación falla, primero clean}
Esos sucesos paranormales de “ayer funcionaba y hoy no” se deben en seis o siete de cada diez casos a una caché de compilación desordenada. Ejecuta ./gradlew clean y vuelve a compilar: cura la mayoría de los misterios. Si persiste, reinicia Android Studio; este remedio funciona mejor combinado con File Invalidate Caches.

\end{tcbinfo}

\booksection{Palabras finales: de la primera App a la App número N}

\booksection{Palabras finales: de la primera App a la App número N}

Si miras atrás: has pasado de “no saber qué es un Activity” a tener el entorno instalado, Kotlin aprendido, interfaces dibujadas con Compose, base de datos y red conectadas a tu App, y encima pruebas, empaquetado y publicación. En este camino ya recorriste el tramo más difícil: el de cero a uno.

El mundo tecnológico cambia rápido: los frameworks se actualizan, las herramientas se renuevan y la IA irá asumiendo cada vez más trabajo de “teclear”. Pero hay algo que no pasa de moda: tu capacidad de descomponer una idea difusa en pasos claros e implementarla paso a paso. Las herramientas son el amplificador; el pensamiento, los cimientos. El amplificador se puede reemplazar; los cimientos no se mueven.

Ya no eres solo "alguien que sabe usar el móvil": puedes convertir las ideas de tu cabeza en apps que funcionan de verdad en trescientos millones de dispositivos Android. Desde hoy, crea con el teclado: ese es el romanticismo de los desarrolladores Android.

\booksection{Referencias}

\begin{enumerate}
  \item Android Developers. (2026). Android Developer Guides and Jetpack Compose Documentation. developer.android.com.
  \item Kotlin Foundation. (2026). Kotlin Language Documentation. kotlinlang.org.
  \item Android Open Source Project. (2026). Android Platform Version History and API Levels.
  \item Google Play. (2026). Play Console Help: Publishing and App Signing.
  \item AndroidX Release Notes. (2026). Room / DataStore / Navigation / WorkManager releases.
  \item Square. (2026). Retrofit Documentation. github.com/square/retrofit.
\end{enumerate}
